\documentclass[10pt,a4paper]{article}
\usepackage{amsmath,amssymb,amsthm,graphicx}
\usepackage{hyperref}
\usepackage{fontspec}
\title{XOR composition lower bounds for an AND of small De Morgan formulas}
\author{Ivan Mihajlin, Anastasia Sofronova}
\date{doi:10.4230/LIPIcs.CCC.2022.13; CC BY 4.0}
\begin{document}
\maketitle
\noindent\textit{Editorial scope.} Counted original Theorem7 (Main lemma) once, retaining both original random-f and existential-g quantifiers and top-AND restriction. Complete well-mixed existence Theorem20 and every source-local supporting proof are included, not separately counted. AND/OR variants and the modified-Andreev consequence are not separate tasks. Original official author mathematics, including its notation, remains unchanged.\par\medskip
\noindent\textit{Allowed standard prerequisites (editorial).}\begin{itemize}
\item Classical counting lower bound for a random Boolean function and multiplicative Chernoff bounds. The author imports counting fact [16]; exact Chernoff statement and every source-local matrix, well-mixing, independence/counting calculation retained.
\end{itemize}\medskip
% Faithful official CCC2022 article13 transcription: Section2 and Sections3.2–4.
% Original author numbering and formulas are retained. No new proof or corrected author claim.
\section*{2 Preliminaries}
\subsection*{2.1 Notation}
Let us recall the definition of the XOR-composition.
\paragraph{Definition 6 ([14]).} XOR-composition $f\boxplus g:\{0,1\}^{2n}\to\{0,1\}$ of two function $f:\{0,1\}^n\to\{0,1\}$ and $g:\{0,1\}^n\to\{0,1\}^n$ is defined as follows:
\[(f\boxplus g)(x,y):=f(g(x)\oplus y).\]
Let us list some notation in regard to formula complexity.
\paragraph{Definition 8.} Let $L(f)$ be the minimum number of leaves in a formula $F$ over basis $\{\wedge,\vee,\neg\}$ such that it computes $f$.
Let $h(x,y)$ be a function of two variables. We denote as $h^x$ a function: $h^x(y):=h(x,y)$.
We also introduce the following shortcut notation for dealing with matrices. Let $M$ be a matrix with rows indexed by $X$ and columns indexed by $Y$. We denote a submatrix with rows indexed by $A$ for $A\subseteq X$ and columns indexed by $Y$ as $M^A$. An element of a matrix, located in row indexed by $x\in X$ and column indexed by $y\in Y$, is denoted as $M[x,y]$. Analogously, we denote a row indexed by $x$ as $M[x]$.
For the rest of the paper, we consider only boolean matrices whose rows and columns are indexed by $X:=\{0,1\}^n$ and $Y:=\{0,1\}^n$.
\paragraph{Definition 9.} For a function $h:\{0,1\}^n\times\{0,1\}^n\to\{0,1\}$, we define a matrix $M_h$:
\[M_h[x,y]:=h(x,y).\]
\paragraph{Definition 10.} For a pair of functions $f:\{0,1\}^n\to\{0,1\}$ and $g:\{0,1\}^n\to\{0,1\}^n$ we define a matrix $M_{f,g}$:
\[M_{f,g}[x,y]:=f(g(x)\oplus y).\]
\paragraph{Definition 11.} For a function $f:\{0,1\}^n\to\{0,1\}$ and a set of functions $Z$ from $\{0,1\}^n\to\{0,1\}^n$ we define a matrix $M_{f,Z}$:
\[M_{f,Z}[x,y]:=\bigvee_{g\in Z}f(g(x)\oplus y).\]
For the rest of the paper, $N:=2^n$.
As we use $f$ and $g$ solely to denote first and second arguments of XOR-composition, we always imply the following domains and ranges for them: $f:\{0,1\}^n\to\{0,1\}$, $g:\{0,1\}^n\to\{0,1\}^n$. Analogously, $x$ and $y$ are implied to be vectors from $\{0,1\}^n$.
\subsection*{2.2 Warm-up lemmas}
\paragraph{Lemma 12.} For a random $f$ and arbitrarily fixed $g$ and $x$ we have: $L((f\boxplus g)^x)\geq N^{1-o(1)}$ with probability $1-o(1)$.
\paragraph{Proof.} As we fix $g$ and $x$, $f(g(x)\oplus y)$ depends only on $y$, so let $h(y):=f(g(x)\oplus y)$.
Any formula for $h$ can be transformed into a formula for $f$ by adding negations to the variables $y_i$ for all $i$ where $g(x)_i=1$. But, as a random function, $f$ does not have a formula of size less than $N^{1-o(1)}$ with high probability [16].\hfill$\triangleleft$
The next lemma gives us connection between formula complexity of a function $f$ and the size of $f^{-1}(0)$:
\paragraph{Lemma 13.} For $f:\{0,1\}^n\to\{0,1\}$, it holds that $L(f)\leq|f^{-1}(0)|\cdot n$.
\paragraph{Proof.} Let us construct a CNF formula for $f$. For any $x\in f^{-1}(0)$, we write a clause of $n$ variables which becomes violated iff we substitute $x$ to those variables. It is easy to check that a CNF formula composed exactly of $|f^{-1}(0)|$ such clauses represents function $f$. This formula has $|f^{-1}(0)|\cdot n$ leaves, which proves the inequality.\hfill$\triangleleft$
\paragraph{Lemma 14.} Let $h(x,y)$ be a function of two variables. Then $L(h)\geq L(h^x)$ for any $x$.
\paragraph{Proof.} Taking a formula, computing $h$, we get a formula, computing $h^x$, by hardwiring the value of $x$ into it.\hfill$\triangleleft$
\section*{3 Proof of the main theorem}
\subsection*{3.2 Proof of Theorem 7}
Now we prove Theorem 7.

\paragraph{Theorem 7 (Main lemma).} For any $0<\alpha<1/5$ and any constant $0<\varepsilon<\alpha$, with probability $1-o(1)$ for a random function $f:\{0,1\}^n\to\{0,1\}$, there exists a function $g:\{0,1\}^n\to\{0,1\}^n$, such that $f\boxplus g$ is not computable by an AND of $2^{(2\alpha-\varepsilon)n}$ formulas of size at most $2^{(1-\alpha-\varepsilon)n}$.

To achieve that, we need to define a notion of well-mixed set of functions.
\paragraph{Definition 19 (Well-mixed set of functions).} A set of functions $G$ from $\{0,1\}^n\to\{0,1\}^n$ is $(Q,D,P)$-well-mixed for $f$ if $\forall Z\subset G$, $|Z|=Q$, there exists a set $K\subset\{0,1\}^n$, $|K|\leq P$, such that $M_{f,Z}^{X\setminus K}$ has no more than $D$ zeroes in total. We call the set $K$ unlucky rows, and we call the set $X\setminus K$ lucky rows.
Informally, we say that a set is well-mixed if all of its subsets behave close to how a random subset of functions would behave.
The key assumption for proving Theorem 7 is that there exists large enough well-mixed set of functions $G$ with suitable parameters.
\paragraph{Theorem 20.} For $0<\alpha<1/5$, let $G$ be a family of functions $\{0,1\}^n\to\{0,1\}^n$ of size $N^{\frac14 N^{1-\alpha}}$, each sampled uniformly at random. Then $G$ is $(N^\alpha,2N^{2-2\alpha},2N^{1-\alpha})$-well-mixed set for a random function $f$ with probability at least $1-\frac1{\exp(\exp(n))}$.
We leave proving this theorem until next section, and for now let us prove Theorem 7 under this assumption.
Let us recall that $N$ stands for $2^n$, as this notation is used heavily below.
\paragraph{Proof.} Let us randomly pick a function $f$. Then we take a set of functions $G$ such that $|G|=N^{\frac14 N^{1-\alpha}}$ and $G$ is $(N^\alpha,2N^{2-2\alpha},2N^{1-\alpha})$-well-mixed. We aim to show that $f\boxplus g$ is hard for some $g\in G$.
We assume the contrary. Suppose for all $g\in G$ the XOR-composition $f\boxplus g$ can be represented as AND of small enough formulas. Formally, for any $g\in G$:
\[f\boxplus g=\bigwedge_{i=1}^{N^{2\alpha-\varepsilon}}h_{g,i}\]
where all $h_{g,i}$ are computable by formulas of size $N^{1-\alpha-\varepsilon}$. For any $g$ we fix the smallest formula for $f\boxplus g$ of such form, and we denote it $F_g$.
We define a measure $C(h)$ of any function $h$ of two arguments:
\[C(h)=\sum_{x\in X}L(h^x).\]
Let us note that $C(f\boxplus g)\geq N\cdot L(f)\geq N^{2-o(1)}$. We prove that this measure is subadditive in the following sense:
\paragraph{Lemma 21.} Let $h(x,y)=\circ(g_1,g_2,\ldots,g_k)(x,y)$, where $\circ$ is either $\wedge$ or $\vee$. Then $C(h)\leq C(g_1)+\ldots+C(g_k)$.
\paragraph{Proof.} If we prove that the inequality holds for any specific $x$, namely, $L(h^x)\leq L(g_1^x)+\ldots+L(g_k^x)$, the lemma immediately follows.
Since $h(x,y)=\circ(g_1,g_2,\ldots,g_k)(x,y)$, we can construct a formula computing $h^x$, applying operation $\circ$ to formulas computing each $g_i^x$. The inequality follows.\hfill$\triangleleft$
It means that for any $g\in G$ and $F_g=\bigwedge_{i=1}^{N^{2\alpha-\varepsilon}}h_{g,i}$ we can fix $1\leq i_g\leq N^{2\alpha-\varepsilon}$ such that for $h_{g,i_g}$ the measure is big enough, namely:
\[C(h_{g,i_g})\geq N^{2-2\alpha+\varepsilon-o(1)}.\]
We are going to arrive at a contradiction by showing that this measure is smaller than it should be. For that, we are going to use two different upper bounds on $L(h_{g,i_g}^x)$ for any $x$. As we assumed that formulas for $h_{g,i_g}$ are small enough, $L(h_{g,i_g})\leq N^{1-\alpha}$. Then for any $x\in X$:
\[L(h_{g,i_g}^x)\leq L(h_{g,i_g})\leq N^{1-\alpha}.\]
On the other hand,
\[L(h_{g,i_g}^x)\leq|(h_{g,i_g}^x)^{-1}(0)|\cdot n\]
by Lemma 13.
Now let $Z_h$ be a subset of $G$ such that $\forall g\in Z_h:h_{g,i_g}=h$ for some fixed $h$.
We consider two cases depending on the size of $Z_h$. First, suppose that $|Z_h|\geq N^\alpha$. We are going to arrive at a contradiction with this assumption. Without losing generality, assume $|Z_h|=N^\alpha$, as we can take a subset of $Z_h$ of exactly this size. Since $G$ is $(N^\alpha,2N^{2-2\alpha},2N^{1-\alpha})$-well-mixed, we can consider a matrix $M_{f,Z_h}$ and distinguish set of unlucky rows $K$ and set of lucky rows $X\setminus K$ in it.
From the properties of well-mixed set, we know that:
\begin{itemize}
\item $|K|\leq2N^{1-\alpha}$;
\item there is an upper bound on number of zeroes in $M_{f,Z_h}^{X\setminus K}$, which are exactly common zeroes of all $g\in Z_h$ on $X\setminus K$:
\[\sum_{x\notin K}\left|\bigcap_{g\in Z_h}(g^x)^{-1}(0)\right|\leq2N^{2-2\alpha}.\]
\end{itemize}
Let us also note that $M_h\geq M_{f,Z_h}$, since $(f\boxplus g)(x,y)=1$ for any $g\in Z_h$ implies $h(x,y)=1$.
\begin{align*}
C(h)&=\sum_xL(h^x)=\sum_{x\in K}L(h^x)+\sum_{x\notin K}L(h^x)\\
&\leq\sum_{x\in K}L(h^x)+\sum_{x\notin K}|(h^x)^{-1}(0)|\cdot n\\
&\leq2N^{1-\alpha}\cdot N^{1-\alpha}+\sum_{x\notin K}\left|\bigcap_{g\in Z_h}(g^x)^{-1}(0)\right|\cdot n\\
&\leq2N^{2-2\alpha}+2N^{2-2\alpha}\cdot n=N^{2-2\alpha+o(1)}
\end{align*}
since we have no more than $2N^{2-2\alpha}$ zeroes in all lucky rows in $M_{f,Z_h}$, therefore, the number of zeroes in rows $x\notin K$ in $M_h$ does not exceed this as well. We know that $C(h)\geq N^{2-2\alpha+\varepsilon-o(1)}$ though, so this is a contradiction.
So $|Z_h|<N^\alpha$, and then the number of different $h$'s is at least $|G|/N^\alpha\geq N^{\frac14N^{1-\alpha}(1-o(1))}$.
Now let $s:=\max_hL(h)$. The number of functions with formulas of size at most $s$ is at most $s^{s(1+o(1))}$, so $s^{s(1+o(1))}\geq N^{\frac14N^{1-\alpha}(1-o(1))}$ and $s\log s(1+o(1))\geq\frac14N^{1-\alpha}\log N(1-o(1))$. Then $s>N^{1-\alpha-\varepsilon}$, and this completes the proof.\hfill$\triangleleft$
\section*{4 The existence of a well-mixed set of functions}
Let us restate the definition of a well-mixed set:
\paragraph{Definition 19 (Well-mixed set of functions).} A set of functions $G$ from $\{0,1\}^n\to\{0,1\}^n$ is $(Q,D,P)$-well-mixed for $f$ if $\forall Z\subset G$, $|Z|=Q$, there exists a set $K\subset\{0,1\}^n$, $|K|\leq P$, such that $M_{f,Z}^{X\setminus K}$ has no more than $D$ zeroes in total. We call the set $K$ unlucky rows, and we call the set $X\setminus K$ lucky rows.
The goal of this section is to prove the following theorem:
\paragraph{Theorem 20.} For $0<\alpha<1/5$, let $G$ be a family of functions $\{0,1\}^n\to\{0,1\}^n$ of size $N^{\frac14N^{1-\alpha}}$, each sampled uniformly at random. Then $G$ is $(N^\alpha,2N^{2-2\alpha},2N^{1-\alpha})$-well-mixed set for a random function $f$ with probability at least $1-\frac1{\exp(\exp(n))}$.
Note that we sample $G$ as a family of functions, with possible repetitions, which we use heavily in the proof. But with probability at least
\[1-\left(\frac{N^{\frac14N^{1-\alpha}}}{N^N}\right)^{N^{\frac14N^{1-\alpha}}}\geq1-\left(\frac1{N^{\frac34N}}\right)^{N^{\frac14N}}\geq e^{-1/N^{N/2}}\geq1-\frac1{N^{N/2}}\]
all functions turn out to be different from each other.
The plan of the proof is as follows:
\begin{itemize}
\item We identify which rows are supposed to be unlucky, and there will be two kinds of those: good and bad.
\item For each kind of unlucky rows, we do three steps:
\begin{itemize}
\item we prove that for a fixed row the probability to be of this kind for a random family of functions is very low;
\item we amplify this probability further, calculating the probability that some number of rows fail us simultaneously;
\item we sum the error over subsets of the random family.
\end{itemize}
\end{itemize}
The first kind of unlucky rows would be a bad row.
\paragraph{Definition 22.} Let $Z$ be a family of functions $\{0,1\}^n\to\{0,1\}^n$, each function sampled uniformly at random. We call $x$ a bad row for $Z$, if $\dim(\operatorname{span}(\bigcup_{g\in Z}\{g(x)\}))\leq2\alpha n$, and a good row otherwise. Here the vector space is defined over $\mathbb F_2$ in a natural way.
We bound the probability that a given $x$ is bad for a random $Z$:
\paragraph{Lemma 23.} Let $Z$ be a random family of functions $\{0,1\}^n\to\{0,1\}^n$ of size $Q$. Then for a fixed $x$:
\[\Pr_Z\left[\dim\left(\operatorname{span}\left(\bigcup_{g\in Z}\{g(x)\}\right)\right)\leq2\alpha n\right]\leq N^{2\alpha n}\left(\frac1{N^{1-2\alpha}}\right)^Q.\]
\paragraph{Proof.} The proof is a simple counting argument. To have a vector subspace of dimension $2\alpha n$, we need to pick $2\alpha n$ generating vectors. There are $2^{2\alpha n}=N^{2\alpha}$ vectors in such space, and for each of $Q$ functions from $Z$ we pick its value in point $x$ from those possibilities, versus $N^Q$ possibilities in general case. So we have:
\begin{align*}
\Pr_Z\left[\dim\left(\operatorname{span}\left(\bigcup_{g\in Z}\{g(x)\}\right)\right)\leq2\alpha n\right]
&\leq\binom N{2\alpha n}(N^{2\alpha})^QN^{-Q}\\
&\leq N^{2\alpha n}\left(\frac1{N^{1-2\alpha}}\right)^Q.
\end{align*}\hfill$\triangleleft$
Now let us prove that with high probability, we have no more than $N^{1-\alpha}$ bad rows for a random $Z$.
\paragraph{Lemma 24.} Let $Z$ be as in Lemma 23. Then
\[\Pr_Z\left[\exists V\subset X,|V|=N^{1-\alpha}:\forall x\in V:\dim\left(\operatorname{span}\bigcup_{g\in Z}\{g(x)\}\right)\leq2\alpha n\right]\leq N^{-(1-2\alpha)QN^{1-\alpha}(1-o(1))}\]
if $Q=\omega(n)$.
\paragraph{Proof.} Since each function in $Z$ is sampled uniformly at random, for each $x$ the event of ``being bad'' is independent of others. So for random $Z$ the probability that a fixed set of $N^{1-\alpha}$ $x$'s is bad can be bounded by $(N^{2\alpha n}(1/N^{1-2\alpha})^Q)^{N^{1-\alpha}}$. After that we apply a union bound over all sets of $x$'s.
We get:
\begin{align*}
N^{N^{1-\alpha}}N^{2\alpha nN^{1-\alpha}}\left(\frac1{N^{1-2\alpha}}\right)^{QN^{1-\alpha}}
&=\left(\frac{N^{2\alpha n+1}}{N^{(1-2\alpha)Q}}\right)^{N^{1-\alpha}}\\
&=N^{-(1-2\alpha)Q(1-\frac{2\alpha n+1}{(1-2\alpha)Q})N^{1-\alpha}}\\
&\leq N^{-(1-2\alpha)QN^{1-\alpha}(1-o(1))}
\end{align*}
for $Q=\omega(n)$.\hfill$\triangleleft$
At last, we apply union bound over all choices of $Z$ from $G$.
\paragraph{Lemma 25.} Let $G$ be as in Theorem 20. Then:
\begin{multline*}
\Pr_G[\exists Z\subset G,|Z|=Q,\exists V\subset X,|V|=N^{1-\alpha}:\forall x\in V:x\text{ is bad for }Z]\\
\leq N^{-\frac Q4N^{1-\alpha}(1-o(1))}.
\end{multline*}
\paragraph{Proof.} For any specific $Z\subset G$ we can apply Lemma 24, as each function in $G$ is sampled uniformly at random. We sum the error over all possible choices of $Z$ and get:
\begin{align*}
|G|^QN^{-(1-2\alpha)QN^{1-\alpha}(1-o(1))}
&=\left(\frac{N^{\frac14N^{1-\alpha}}}{N^{(1-2\alpha)N^{1-\alpha}(1-o(1))}}\right)^Q\\
&\leq N^{-\frac Q4N^{1-\alpha}(1-o(1))}.
\end{align*}\hfill$\triangleleft$
Hence, for every $Z\subset G$ there are no more than $N^{1-\alpha}$ bad rows with high probability.
We consider bad rows unlucky. After we take those out of consideration, we prove that almost all good rows are lucky.
To do that, first we consider a technical Lemma:
\paragraph{Lemma 26.} Let $f$ be a function $\{0,1\}^n\to\{0,1\}$, chosen uniformly randomly, $S$ be any set of functions from $\{0,1\}^n\to\{0,1\}^n$ such that for a fixed $x$ values $g(x)$ for all $g\in S$ are linearly independent, $2\leq|S|<n/2$. Then:
\[\Pr_f[M_{f,S}[x]\text{ has }2N^{2-|S|}\text{ zeroes}]\leq2^{(n+1)|S|-2^{n-2|S|-2}}.\]
This is a corollary from the following statement by Kaave Hosseini (from personal communication):
\paragraph{Lemma 27.} Let $B\subset\{0,1\}^n$ be a random set, where each point lies with probability $1/2$, $L=x_1,\ldots,x_k$ be linearly independent vectors, $k\geq2$. Then:
\[\Pr_B(|(B+x_1)\cap\ldots\cap(B+x_k)|\geq2^{n-k+1})\leq2^{(n+1)k-2^{n-2k-2}}.\]
First we show how Lemma 26 follows from Lemma 27.
\paragraph{Proof of 26 from 27.} The value of $M_{f,S}[x,y]$ is zero iff $\forall g\in S:f(g(x)\oplus y)=0$. Plugging $B:=f^{-1}(0)$, $k:=|S|$, $L:=\{g(x)\}_{g\in S}$ in 27, we get 26, since zeroes of $f(g(x)\oplus y)$ are shifted by vector $g(x)$ in comparison with zeroes of $f$.
Now we prove Lemma 27.
\paragraph{Proof.} To upper bound $|(B+x_1)\cap\ldots\cap(B+x_k)|$, let us consider the following sum:
\[A:=\sum_{y\in\{0,1\}^n}\chi(y+x_1)\chi(y+x_2)\ldots\chi(y+x_k)\]
where $\chi$ is a characteristic function of a set $B$. This sum equals $|(B+x_1)\cap\ldots\cap(B+x_k)|$ exactly.
Let us now extend set $L=\{x_1,\ldots,x_k\}$ to a basis and split $\{0,1\}^n$ onto $2^k$ parts, depending on whether a vector contains $x_i$ in its decomposition onto basis vectors or not. Each part contains $2^{n-k}$ different points. Let us consider one of those parts, without losing generality we pick:
\[P=\{y\mid y\text{ does not contain any of }x_i\text{ in decomposition}\}.\]
In the sum $A_P:=\sum_{y\in P}\chi(y+x_1)\ldots\chi(y+x_k)$ every point in $\{0,1\}^n$ occurs as an argument of $\chi$ no more than once, so all summands and all multipliers in them are independent.
If we interpret every summand as a Bernoulli variable which equals $1$ with probability $1/2^k$, we can apply Chernoff bounds.
The exact form that we use here is the following:
\paragraph{Theorem 28 (Chernoff bounds, multiplicative form).} For independent random $X_1,\ldots,X_n$ from $\{0,1\}$ and $0\leq\delta\leq1$:
\[\Pr\left[\sum_iX_i\geq(1+\delta)\mu\right]\leq e^{-\delta^2\mu/3}\]
where $\mu$ is the expected value of $\sum_iX_i$.
Here the expected value of $A_P$ equals $2^{n-2k}$.
\[\Pr_B(A_P>(1+\delta)2^{n-2k})<e^{(-\delta^2)2^{n-2k-2}}.\]
This holds regardless of the choice of a part $P$. We apply union bound and sum the error over all possible parts:
\begin{align*}
\Pr_B(A>(1+\delta)2^{n-k})&\leq\Pr_B(\exists P:A_P>(1+\delta)2^{n-2k})\\
&<2^k\cdot e^{(-\delta^2)2^{n-2k-2}}.
\end{align*}
At last, we apply union bound over all possible choices of $L$:
\[\Pr_B(\exists L:A>(1+\delta)2^{n-k})<2^{nk}\cdot2^k\cdot e^{(-\delta^2)2^{n-2k-2}}.\]
We pick $\delta:=1$ and get:
\[\Pr_B(A>2^{n-k+1})<2^{(n+1)k-2^{n-2k-2}}.\]
This finishes the proof of the Lemma.\hfill$\triangleleft$
We use the Lemma to bound the number of zeroes in good rows. First, let us remember that $\alpha$ is separated from $1/5$ by some small constant. Let that constant be $\gamma$, so $\alpha+\gamma=1/5$.
\paragraph{Lemma 29.} Let $x$ be a good row for a family of functions $Z$. Then:
\[\Pr_f(M_{f,Z}[x]\text{ has }2N^{1-2\alpha}\text{ zeroes})\leq2^{-N^{\alpha+5\gamma}(1-o(1))}.\]
\paragraph{Proof.} Since $x$ is good for $Z$, we can find a subset $S\subset Z$ such that $|S|=2\alpha n$ and $\{g(x)\}_{g\in S}$ are linearly independent. By Lemma 26 and the fact that $n-2\cdot2\alpha n\geq(\alpha+5\gamma)n$, we immediately have an upper bound on number of zeroes in a row $x$:
\begin{align*}
\Pr_f(M_{f,Z}[x]\text{ has }2N^{1-2\alpha}\text{ zeroes})
&\leq\Pr_f(M_{f,S}[x]\text{ has }2N^{1-2\alpha}\text{ zeroes})\\
&\leq2^{(n+1)2\alpha n-2^{(\alpha+5\gamma)n-2}}\\
&=2^{-N^{\alpha+5\gamma}(1-o(1))}.
\end{align*}\hfill$\triangleleft$
\paragraph{Lemma 30.} Let $Z$ be a random family of functions from $\{0,1\}^n\to\{0,1\}^n$. Then:
\begin{multline*}
\Pr_{f,Z}[\exists V,V\text{ is a set of good rows for }Z,|V|=N^{1-\alpha}:\\
\forall x\in V:M_{f,Z}[x]\text{ has }2N^{1-2\alpha}\text{ zeroes}]\leq2^{-N^{1+5\gamma}(1-o(1))}.
\end{multline*}
\paragraph{Proof.} Again, as $Z$ is picked uniformly at random, the events regarding every $x$ are independent for a fixed $V$. Now let us sum the error over all choices of $V$:
\begin{multline*}
\Pr_{f,Z}[\exists V,V\text{ is a set of good rows for }Z,|V|=N^{1-\alpha}:\\
\forall x\in V:M_{f,Z}[x]\text{ has }2N^{1-2\alpha}\text{ zeroes}]\\
\leq\binom N{N^{1-\alpha}}\left(2^{-N^{\alpha+5\gamma}(1-o(1))}\right)^{N^{1-\alpha}}\\
\leq\binom N{N^{1-\alpha}}2^{-N^{1+5\gamma}(1-o(1))}\\
\leq2^{N^{1-\alpha}\cdot\log N}\cdot2^{-(1-o(1))N^{1+5\gamma}}\leq2^{-(1-o(1))N^{1+5\gamma}}.
\end{multline*}\hfill$\triangleleft$
Now, we are ready to sum over all possible choices of $Z$ in a random $G$.
\paragraph{Lemma 31.} Let $G$ be as in Theorem 20. Then:
\begin{multline*}
\Pr_f[\exists Z\subset G,|Z|=Q,\exists V,V\text{ is a set of good rows for }Z,|V|=N^{1-\alpha}:\\
\forall x\in V:M_{f,Z}[x]\text{ has }2N^{1-2\alpha}\text{ zeroes}]\leq N^{Q\frac14N^{1-\alpha}}2^{-(1-o(1))N^{1+5\gamma}}.
\end{multline*}
\paragraph{Proof.} This is a union bound over all choices of a subset $Z$ of size $Q$ from set $G$ of size $N^{\frac14N^{1-\alpha}}$.\hfill$\triangleleft$
Now we are ready to prove Theorem 20.
Let us take a family $G$ of functions from $\{0,1\}^n\to\{0,1\}^n$ of size $N^{\frac14N^{1-\alpha}}$, where each function is sampled uniformly at random. With probability $1-1/\exp(\exp(n))$ all functions turn out to be different from each other.
Let us take any $Z\subset G$ of size $N^\alpha$. We plug $Q:=N^\alpha$ into Lemma 25 and get that with probability at least $1-N^{-N^{1-\alpha}}$ we have no more than $N^{1-\alpha}$ bad rows for our $Z$.
Plugging this parameter in Lemma 31, we get that with probability $1-N^{\frac14N-N^{1+5\gamma}(1-o(1))}=1-2^{-N^{1+5\gamma}(1-o(1))}$ no more than $N^{1-\alpha}$ good rows have $2N^{1-2\alpha}$ zeroes in them.
So, with very high probability, all properties hold. We say that both bad rows and good rows with at least $2N^{1-2\alpha}$ zeroes are our unlucky rows, and this is a set $K$ from Theorem 20.
Now, as every lucky row has at most $2N^{1-2\alpha}$ zeroes, and there are no more than $N$ lucky rows, there are no more than $2N^{2-2\alpha}$ zeroes in all lucky rows, which gives us the statement of Theorem 20.

\begin{thebibliography}{18}
\bibitem[1]{ref1} Alexander E. Andreev. On a method for obtaining more than quadratic effective lower bounds for the complexity of $\pi$-schemes. Moscow University Mathematics Bulletin, 42(1):24–29, 1987.
\bibitem[2]{ref2} Susanna F. de Rezende, Or Meir, Jakob Nordström, Toniann Pitassi, and Robert Robere. KRW composition theorems via lifting. In Sandy Irani, editor, 61st IEEE Annual Symposium on Foundations of Computer Science, FOCS 2020, Durham, NC, USA, November 16-19, 2020, pages 43–49. IEEE, 2020. doi:10.1109/FOCS46700.2020.00013.
\bibitem[3]{ref3} Irit Dinur and Or Meir. Toward the KRW composition conjecture: Cubic formula lower bounds via communication complexity. In Ran Raz, editor, 31st Conference on Computational Complexity, CCC 2016, May 29 to June 1, 2016, Tokyo, Japan, volume 50 of LIPIcs, pages 3:1–3:51. Schloss Dagstuhl - Leibniz-Zentrum für Informatik, 2016. doi:10.4230/LIPIcs.CCC.2016.3.
\bibitem[4]{ref4} Jeff Edmonds, Russell Impagliazzo, Steven Rudich, and Jirí Sgall. Communication complexity towards lower bounds on circuit depth. Comput. Complex., 10(3):210–246, 2001. doi:10.1007/s00037-001-8195-x.
\bibitem[5]{ref5} Dmitry Gavinsky, Or Meir, Omri Weinstein, and Avi Wigderson. Toward better formula lower bounds: The composition of a function and a universal relation. SIAM J. Comput., 46(1):114–131, 2017. doi:10.1137/15M1018319.
\bibitem[6]{ref6} Johan Håstad. The shrinkage exponent of de morgan formulas is 2. SIAM J. Comput., 27(1):48–64, 1998. doi:10.1137/S0097539794261556.
\bibitem[7]{ref7} Johan Håstad and Avi Wigderson. Composition of the universal relation. In Advances In Computational Complexity Theory, Proceedings of a DIMACS Workshop, New Jersey, USA, December 3-7, 1990, pages 119–134, 1990. doi:10.1090/dimacs/013/07.
\bibitem[8]{ref8} Russell Impagliazzo, Raghu Meka, and David Zuckerman. Pseudorandomness from shrinkage. In 2012 IEEE 53rd Annual Symposium on Foundations of Computer Science, pages 111–119. IEEE, 2012.
\bibitem[9]{ref9} Russell Impagliazzo and Noam Nisan. The effect of random restrictions on formula size. Random Struct. Algorithms, 4:121–134, 1993.
\bibitem[10]{ref10} Mauricio Karchmer, Ran Raz, and Avi Wigderson. Super-logarithmic depth lower bounds via the direct sum in communication complexity. Computational Complexity, 5(3/4):191–204, 1995. doi:10.1007/BF01206317.
\bibitem[11]{ref11} Mauricio Karchmer and Avi Wigderson. Monotone circuits for connectivity require super- logarithmic depth. In Proceedings of the 20th Annual ACM Symposium on Theory of Computing, May 2-4, 1988, Chicago, Illinois, USA, pages 539–550, 1988. doi:10.1145/62212.62265.
\bibitem[12]{ref12} Valeriy Mihailovich Khrapchenko. Complexity of the realization of a linear function in the class of II-circuits. Mathematical Notes of the Academy of Sciences of the USSR, 9(1):21–23, 1971.
\bibitem[13]{ref13} Sajin Koroth and Or Meir. Improved composition theorems for functions and relations. In Eric Blais, Klaus Jansen, José D. P. Rolim, and David Steurer, editors, Approximation, Randomiza- tion, and Combinatorial Optimization. Algorithms and Techniques, APPROX/RANDOM 2018, August 20-22, 2018 - Princeton, NJ, USA, volume 116 of LIPIcs, pages 48:1–48:18. Schloss Dag- stuhl - Leibniz-Zentrum für Informatik, 2018. doi:10.4230/LIPIcs.APPROX-RANDOM.2018.48.
\bibitem[14]{ref14} Ivan Mihajlin and Alexander Smal. Toward better depth lower bounds: The XOR-KRW conjecture. In Valentine Kabanets, editor, 36th Computational Complexity Conference, CCC 2021, July 20-23, 2021, Toronto, Ontario, Canada (Virtual Conference), volume 200 of LIPIcs, pages 38:1–38:24. Schloss Dagstuhl - Leibniz-Zentrum für Informatik, 2021. doi:10.4230/LIPIcs.CCC.2021.38.
\bibitem[15]{ref15} Michael S. Paterson and Uri Zwick. Shrinkage of de morgan formulae under restriction. In Proceedings of the 32nd Annual Symposium on Foundations of Computer Science, SFCS ’91, pages 324–333, USA, 1991. IEEE Computer Society. doi:10.1109/SFCS.1991.185385.
\bibitem[16]{ref16} J. Riordan and C. Shannon. The number of two-terminal series-parallel networks. J Math Phys, 21, January 1942. doi:10.1002/sapm194221183.
\bibitem[17]{ref17} Bella Abramovna Subbotovskaya. Realization of linear functions by formulas using $\wedge$, $\vee$, $\neg$. In Doklady Akademii Nauk, volume 136-3, pages 553–555. Russian Academy of Sciences, 1961.
\bibitem[18]{ref18} Avishay Tal. Shrinkage of de morgan formulae by spectral techniques. Proceedings - Annual IEEE Symposium on Foundations of Computer Science, FOCS, pages 551–560, December 2014. doi:10.1109/FOCS.2014.65.
\end{thebibliography}
\end{document}
